\chapter{Conclusão}
\label{chap7}

\section{Conclusão}

O desenvolvimento deste trabalho permitiu a construção de um sistema web voltado ao apoio da gestão do Programa de Aquisição de Alimentos na modalidade Doação Simultânea, tendo como referência o contexto do Assentamento PDS Osvaldo de Oliveira. A proposta partiu de uma necessidade concreta: reduzir a dependência de controles manuais, centralizar informações da chamada pública e facilitar o acompanhamento das entregas realizadas pelas famílias agricultoras aos destinos vinculados ao edital. Ao longo do processo, o sistema deixou de ser apenas uma ideia inicial e passou a constituir uma aplicação em fase de homologação, com funcionalidades implementadas e com possibilidade de validação junto aos usuários envolvidos.

Entre os resultados alcançados, destacam-se a implementação de fluxos essenciais para o funcionamento da plataforma, funcionalidades que representam uma base importante para a gestão do PAA, pois permitem organizar informações que antes eram acompanhadas principalmente por planilhas e registros manuais. Do ponto de vista técnico, o trabalho atingiu um resultado relevante ao estruturar a aplicação em módulos separados, com contextos bem definidos. A escolha por uma arquitetura modular favoreceu a independência entre a interface e as regras de negócio, permitindo que novas funcionalidades fossem incorporadas de forma progressiva.

Além da implementação, o projeto evidenciou a importância de aproximar o desenvolvimento de software do contexto social em que a ferramenta será utilizada. O levantamento de requisitos, as reuniões com os gestores do PAA e os ajustes realizados a partir de novas demandas evidenciam que a construção de uma solução não depende apenas da escolha de tecnologias, mas também da compreensão do processo real de trabalho, agregando valor o que foi desenvolvido.

Os trabalhos futuros devem partir da análise do que já foi implementado e do estágio atual do sistema. Para as novas implementações, a metodologia participativa deverá continuar presente, para que o software tenha maior aderência aos usuários e faça diferença no dia a dia dos coordenadores, administradores e agricultores. A análise da arquitetura planejada e implementada indica que ela atende a uma demanda simples. No entanto, em um contexto mais amplo, considerando o Brasil, o sistema deverá implementar métricas de desempenho, reforçar a segurança dos dados, dividir funcionalidades em serviços e contar com uma estrutura capaz de suportar milhares de requisições, caso seja necessário.

Por fim, conclui-se que este trabalho contribui tanto para a organização técnica de uma aplicação web quanto para a construção de uma ferramenta voltada a uma demanda social específica. O sistema desenvolvido oferece uma base para tornar a gestão das chamadas públicas mais transparente, menos dependente de processos manuais e mais próxima das necessidades dos usuários. Ao mesmo tempo, sua continuidade exigirá novas etapas de amadurecimento, especialmente para que a solução possa deixar de ser apenas um sistema em homologação e se tornar uma ferramenta estável.

